\chapter{Lattice Geometry and Inscribed Polygons}
\label{chap:polygon-geometry}

\section{The Inscribed Polygon Universe}

Let $\mathbb{Z}_n = \{0, 1, \dots, n-1\}$ denote the cyclic ring of
integers modulo $n$, drawn as the vertices of a regular
$n$-gon in the Euclidean plane, oriented with vertex $0$ at the top
and indices advancing clockwise. The physical symmetries of this circular
lattice constitute the dihedral group $D_n$ of order $2n$, generated by
the unit rotation $\rho \colon x \mapsto x + 1 \pmod n$ and the vertical
bilateral reflection $\tau \colon x \mapsto -x \equiv n - x \pmod n$.

An \textbf{inscribed polygon} in $\mathbb{Z}_n$ is specified by an
arbitrary subset of vertices $S \subseteq \mathbb{Z}_n$ of cardinality
$k = |S| \in \{0, 1, \dots, n\}$:
\begin{itemize}
    \item For $k = 0$, $S = \emptyset$ represents the empty polygon.
    \item For $k = 1$, $S = \{v\}$ represents a unigon (isolated vertex).
    \item For $k = 2$, $S = \{u, v\}$ represents a digon (a line of circular span $\min(|u-v|, n-|u-v|)$).
    \item For $3 \le k \le n-1$, $S$ represents a convex $k$-gon inscribed within the lattice.
    \item For $k = n$, $S = \mathbb{Z}_n$ spans the complete $n$-gon.
\end{itemize}

\noindent\textbf{Degenerate Base Lattices: $\mathbb{Z}_1$ and $\mathbb{Z}_2$}\\[0.5em]

The degenerate base rings $\mathbb{Z}_1$ and $\mathbb{Z}_2$ admit trivial
dihedral symmetry groups ($D_1 \cong C_2$ and $D_2 \cong C_2 \times
C_2$). While excluded from the main catalogue plates, their complete
subset universes are fully enumerated:

\begin{figure}[htbp]
\centering
\begin{tikzpicture}[scale=0.8, baseline=-0.5ex]
    % Z_1: empty
    \begin{scope}[shift={(0,0)}]
        \draw[gray!40, line width=0.5pt] (0,0) circle (0.45cm);
        \draw[gray!40, line width=0.4pt] (-0.1,-0.1) -- (0.1,0.1);
        \node[font=\footnotesize\sffamily] at (0,-0.75) {$\mathbb{Z}_1\text{-}0\text{-}1: \emptyset$};
    \end{scope}
    % Z_1: unigon
    \begin{scope}[shift={(2.5,0)}]
        \draw[gray!40, line width=0.5pt] (0,0) circle (0.45cm);
        \fill[mosaicblue!85!black] (0,0.45) circle (2.2pt);
        \node[font=\footnotesize\sffamily] at (0,-0.75) {$\mathbb{Z}_1\text{-}1\text{-}1: \{0\}$};
    \end{scope}
    % Divider
    \draw[gray!30, dashed] (4.5,-0.9) -- (4.5,0.9);
    % Z_2: empty
    \begin{scope}[shift={(6.5,0)}]
        \draw[gray!35, line width=0.5pt] (0,0.45) -- (0,-0.45);
        \draw[gray!40, line width=0.4pt] (-0.1,-0.1) -- (0.1,0.1);
        \node[font=\footnotesize\sffamily] at (0,-0.75) {$\mathbb{Z}_2\text{-}0\text{-}1: \emptyset$};
    \end{scope}
    % Z_2: unigon
    \begin{scope}[shift={(9.0,0)}]
        \draw[gray!35, line width=0.5pt] (0,0.45) -- (0,-0.45);
        \fill[mosaicblue!85!black] (0,0.45) circle (2.2pt);
        \fill[gray!45] (0,-0.45) circle (1.2pt);
        \node[font=\footnotesize\sffamily] at (0,-0.75) {$\mathbb{Z}_2\text{-}1\text{-}1: \{0\}$};
    \end{scope}
    % Z_2: digon
    \begin{scope}[shift={(11.5,0)}]
        \draw[mosaicblue, line width=2.4pt, line cap=round] (0,0.45) -- (0,-0.45);
        \draw[black!80, line width=0.4pt, line cap=round] (0,0.45) -- (0,-0.45);
        \fill[mosaicblue!85!black] (0,0.45) circle (2.0pt);
        \fill[mosaicblue!85!black] (0,-0.45) circle (2.0pt);
        \node[font=\footnotesize\sffamily] at (0,-0.75) {$\mathbb{Z}_2\text{-}2\text{-}1: \{0,1\}$};
    \end{scope}
\end{tikzpicture}
\caption{\label{fig:degenerate-rings}All Inscribed Subsets of $\mathbb{Z}_1$ (2 Orbits) and $\mathbb{Z}_2$ (3 Orbits)}
\end{figure}

\section{The Cauchy--Frobenius Polynomial \texorpdfstring{$CF_n(x)$}{CFn(x)}}

In combinatorial enumeration under group action, the total number of
non-equivalent $k$-element subsets under the dihedral group $D_n$ is
governed by the Cauchy--Frobenius lemma (often historically misattributed
to Burnside) applied via P\'olya's cycle index $Z(D_n; z_1, z_2, \dots, z_n)$.

Evaluating the cycle index under the binary substitution $z_k \mapsto
1 + x^k$ produces the \emph{Cauchy--Frobenius bracelet generating polynomial}:

\begin{equation}
CF_n(x) = Z\left(D_n; 1 + x, 1 + x^2, \dots, 1 + x^n\right) = \sum_{k=0}^n c_{n,k} x^k
\end{equation}
where $c_{n,k}$ is the exact number of dihedral orbits of inscribed $k$-gons.

For odd $n = 2m + 1$, every reflection possesses one fixed vertex and $m$ chords:
\begin{equation}
Z(D_n) = \frac{1}{2n} \left( \sum_{d \mid n} \phi(d) z_d^{n/d} + n z_1 z_2^{(n-1)/2} \right)
\end{equation}
For even $n = 2m$, the $n$ reflections decompose into $n/2$ vertex-vertex axes and $n/2$ edge-edge axes:
\begin{equation}
Z(D_n) = \frac{1}{2n} \left( \sum_{d \mid n} \phi(d) z_d^{n/d} + \frac{n}{2} z_1^2 z_2^{(n-2)/2} + \frac{n}{2} z_2^{n/2} \right)
\end{equation}

Evaluating $CF_n(1)$ by setting $x = 1$ yields the total number of dihedral subset orbits on $\mathbb{Z}_n$.

\section{Complement Duality and Bilateral Symmetry}

Two geometric dualities govern the atlas:

\begin{enumerate}
    \item \textbf{Complement Duality ($k \longleftrightarrow n - k$)}:
    The mapping $S \mapsto \mathbb{Z}_n \setminus S$ defines an isometric bijection between $k$-gons and $(n-k)$-gons. Because the cycle index polynomial is palindromic:
    \begin{equation}
    c_{n,k} \equiv c_{n, n-k}
    \end{equation}
    Every $k$-gon has a canonical complementary $(n-k)$-gon. When $n$ is even, polygons at the central equator $k = n/2$ can be \emph{self-complementary} under dihedral symmetry.
    
    \item \textbf{Bilateral Reflection Symmetry (Chirality)}:
    An inscribed polygon $S$ is \textbf{achiral} (shaded in blue) if there exists a reflection $\tau \in D_n$ such that $\tau(S) = S$. Absent such reflection, $S$ is \textbf{chiral} (shaded in pink), existing as one member of an enantiomorphic pair $(S, \tau(S))$ that split into distinct orbits under cyclic rotations $C_n$.
\end{enumerate}
